\section{Постановка задачи}
\label{sec:Chapter2} \index{Chapter2}

%Необходимо формально изложить суть задачи в данной секции, предоставив такие
%ясные и точные описания, которые позволят в последующем оценить, насколько
%разработанное решение соответствует поставленной задаче. Текст главы должен
%следовать структуре технического задания, включая как описание самой задачи,
%так и набор требований к ее решению.

Пусть $2 \leq Q \leq q$. Определим отображение $\Map: \Sigma_Q \to \Z_q$, как:
\begin{equation*}
	\Map(s) = \left\lfloor\frac{sq}{Q}\right\rfloor
\end{equation*}

В обратную сторону, положим $\Demap: \Z_q \to \Sigma_Q$, как прообраз ближайшего по метрике Ли в $\Z_q$ остатка из $\Map(\Sigma_Q)$. Если таких несколько, выберем минимальный из них (формально: после каноничного $\Z_q \hookrightarrow \Z$).

После этого сообщение $m \in \Sigma_Q^n$ станем кодировать и декодировать посимвольно $\Map: \Sigma_Q^n \to \mathcal{R}_q^n$ и $\Demap: \mathcal{R}_q^n \to \Sigma_Q^n$.

Теперь можно описать протокол для передачи сообщений:

\begin{algorithm}
	\caption{Генерация ключей}
	\begin{algorithmic}
		\Require $q, k \in \N$
		\item $\bold a \leftarrow \Z_q$
		\item $\bold s, \bold e \leftarrow \chi_k$
		\item $\bold b = \bold a\bold s+\bold e$
		\Ensure $pk = (\bold b, \bold a), sk = \bold s$
	\end{algorithmic}
\end{algorithm}

\begin{algorithm}
	\caption{Шифрование}
	\begin{algorithmic}
		\Require $m \in \Sigma_Q^n, \; pk$
		\item $\bold{s'}, \bold{e'}, \bold{e''} \leftarrow \chi_k$
		\item $\bold u = \bold a\bold{s'}+\bold{e'}$
		\item $\bold v = \bold b\bold{s'}+\bold{e''}+\Map(m)$
		\Ensure $c = (\bold u, \bold v)$
	\end{algorithmic}
\end{algorithm}

\begin{algorithm}
	\caption{Разшифрование}
	\begin{algorithmic}
		\Require $c, \; sk$
		\item $\widetilde m = \Demap(\bold v-\bold u\bold s)$
		\Ensure $\widetilde{m}$
	\end{algorithmic}
\end{algorithm}

При этом шумом $\bold z$ назовём $\bold v-\bold u\bold s-\Map(m)$. Имеем:
\begin{equation*}
	\bold z = (\bold b\bold{s'}+\bold{e''})-(\bold a\bold{s'}+\bold{e'})\bold s = \bold e\bold{s'}-\bold{e'}\bold s+\bold{e''}
\end{equation*}

Событие $\widetilde m_i \neq m_i$ обозначим $\mathcal{E}_i$. Заметим, что оно $\eqv \Demap(\bold z)_i~\!\neq~\!0$. Таким образом, можно рассматривать процесс шифрования и расшифрования, как $\mathsf{RLWE}$-канал, который при передаче по нему сообщения $\Map(m)$ из $\mathcal{R}_q^n$ добавляет к нему шум $\bold z$ также из $\mathcal{R}_q^n$.

За сложность взлома протоколов из данного семейства отвечают некоторые хорошие свойства $\mathcal{R}_q^n$, изученные 
